\section{5. Conclusiones y trabajo
futuro}\label{conclusiones-y-trabajo-futuro}

\subsection{5.1 Conclusiones}\label{conclusiones}

\subsection{5.2 Limitaciones de la
solución}\label{limitaciones-de-la-soluciuxf3n}

\begin{itemize}
\item
  \textbf{Dependencia de herramientas externas.} EVMAudit depende del
  correcto funcionamiento de Slither, Mythril y Echidna. Cualquier
  modificación en dichas herramientas puede afectar al comportamiento
  del sistema.
\item
  \textbf{Correlación basada en reglas.} La correlación implementada
  utiliza únicamente:

  \begin{itemize}
  \tightlist
  \item
    contrato
  \item
    función
  \item
    SWC
  \end{itemize}

  Esto puede provocar que determinadas vulnerabilidades relacionadas no
  se agrupen correctamente.
\item
  \textbf{Falsos positivos y falsos negativos.} La herramienta no
  elimina completamente los falsos positivos inherentes a las
  herramientas integradas.
\item
  \textbf{Limitaciones del fuzzing.} Determinadas vulnerabilidades, como
  la reentrancia, requieren contratos atacantes específicos y no pueden
  validarse automáticamente.
\item
  \textbf{Ausencia de análisis inter-contrato.} Actualmente la
  herramienta analiza cada contrato de forma independiente.
\item
  \textbf{Priorización limitada.} La priorización se basa en la
  severidad y en el número de herramientas que detectan la
  vulnerabilidad, sin utilizar métricas más avanzadas como CVSS.
\end{itemize}

\subsection{5.3 Líneas de trabajo
futuro}\label{luxedneas-de-trabajo-futuro}

\begin{itemize}
\tightlist
\item
  incorporar nuevas herramientas
\item
  análisis inter-contrato\\
\item
  CVSS
\item
  Machine Learning
\item
  Semgrep
\item
  Manticore
\item
  dashboard web
\end{itemize}
